\chapter{Demo and GitHub Repository}
\label{app:code}

A demo of the leak identification method is open-source and publicly available at: \\

{ \centering 
\href{https://github.com/AndersMrk/Pipe-Leak-Identification}{\faGithub\ \texttt{github.com/AndersMrk/Pipe-Leak-Identification}} \par
}

\subsection*{How to Run the Demo}

\begin{enumerate}
    \item Clone the repository:
    \begin{lstlisting}[language=bash]
git clone https://github.com/AndersMrk/Pipe-Leak-Identification.git
cd Pipe-Leak-Identification
    \end{lstlisting}

    \item Create a Python virtual environment:
    \begin{lstlisting}[language=bash]
python -m venv .venv
    \end{lstlisting}

    \item Activate the virtual environment.
        \begin{lstlisting}[language=bash]
Windows PowerShell:
    .\.venv\Scripts\Activate.ps1
        \end{lstlisting}
        
        \begin{lstlisting}[language=bash]
Windows Command Prompt:
    .\.venv\Scripts\activate.bat
        \end{lstlisting}
    
        \begin{lstlisting}[language=bash]
Linux/macOS Bash:
    source .venv/bin/activate
        \end{lstlisting}

    \item Install the required dependencies:
    \begin{lstlisting}[language=bash]
python -m pip install -r requirements.txt
    \end{lstlisting}

    \item Run the demo:
    \begin{lstlisting}[language=bash]
python demo.py
    \end{lstlisting}
\end{enumerate}

To test different configurations, these parameters can be changed in \texttt{demo.py}:

\begin{lstlisting}[language=bash]
     TRUE_LEAK_Z = [0.2, 0.6, 0.8]
     TRUE_LEAK_C = [0.02, 0.05, 0.06]
     PIPE_LEN = 1
     THETA = 0.2         
\end{lstlisting}
